\section{Phonons}\label{sec:5.3}

A silicon nucleus in a crystal cooled to a few millikelvin is not at rest. Its displacement along a cube axis has a root-mean-square value of \SI{0.05}{\angstrom} (\cref{we:5-kappa}), about a hundredth of the lattice constant, and at room temperature \SI{0.077}{\angstrom}. The classical normal modes of \cref{sec:5.2} cannot account for the first number, since at zero temperature a classical crystal sits motionless at the bottom of its energy surface. They also miss what matters most for detection: the energy of a mode comes in quanta. A dark-matter particle that scatters in a crystal at zero temperature cannot leave an arbitrary amount of energy in the lattice. It creates phonons, each with the energy $\omega_{\vect{Q}\nu}$ of one mode, out of a lattice in its ground state, and the amplitude for doing so is set by the zero-point motion of that mode. Both facts follow from quantizing the normal modes, which turns the crystal into a set of the harmonic oscillators of \cref{sec:2.3,sec:2.5}.

\subsection{Normal coordinates}

In the harmonic approximation the nuclei move on the quadratic energy surface \cref{eq:5.2-expand}, and the nuclear Hamiltonian of \cref{eq:5.1-BO} becomes
\begin{equation}
  H_{\mathrm{N}} = \sum_{\vect{R}\kappa}\frac{\vect{p}_{\vect{R}\kappa}^2}{2M_\kappa}
  + \frac12\sum_{\vect{R}\kappa\alpha}\sum_{\vect{R}'\kappa'\beta}\Phi^{(2),\alpha\beta}_{\kappa\kappa'}(\vect{R},\vect{R}')\,u^\alpha_{\vect{R}\kappa}u^\beta_{\vect{R}'\kappa'} ,
  \qquad
  [u^\alpha_{\vect{R}\kappa}, p^\beta_{\vect{R}'\kappa'}] = \ii\,\delta_{\vect{R}\vect{R}'}\delta_{\kappa\kappa'}\delta_{\alpha\beta} ,
  \label{eq:5.3-H}
\end{equation}
where $\vect{p}_{\vect{R}\kappa} = -\ii\nabla_{\vect{R}\kappa}$ is the momentum of the nucleus and the constant $E^{\mathrm{ad}}(0)$ has been dropped. The displacements and momenta are now operators. The classical solution of \cref{sec:5.2} already tells us which combinations of them decouple: mass-weighted plane waves in the cell index, projected on the polarization vectors. We define, for each $\vect{Q}$ in the first zone and each branch $\nu$,
\begin{equation}
  u_{\vect{Q}\nu} = \frac{1}{\sqrt N}\sum_{\vect{R}\kappa}\sqrt{M_\kappa}\;\ee^{-\ii\vect{Q}\cdot\vect{R}}\,\vect{e}_{\kappa\nu}(\vect{Q})^*\cdot\vect{u}_{\vect{R}\kappa} , \qquad
  p_{\vect{Q}\nu} = \frac{1}{\sqrt N}\sum_{\vect{R}\kappa}\frac{1}{\sqrt{M_\kappa}}\;\ee^{\ii\vect{Q}\cdot\vect{R}}\,\vect{e}_{\kappa\nu}(\vect{Q})\cdot\vect{p}_{\vect{R}\kappa} ,
  \label{eq:5.3-normal}
\end{equation}
the normal coordinate and the normal momentum of the mode $\vect{Q}\nu$. The transformation is inverted with the lattice sum \cref{eq:3.1-latticesum}, which gives $\sum_{\vect{R}}\ee^{\ii(\vect{Q}' - \vect{Q})\cdot\vect{R}} = N\delta_{\vect{Q}\vect{Q}'}$ for two wavevectors of the first zone, and with the completeness relation of \cref{eq:5.2-ortho}:
\begin{equation}
  \vect{u}_{\vect{R}\kappa} = \frac{1}{\sqrt{NM_\kappa}}\sum_{\vect{Q}\nu}\vect{e}_{\kappa\nu}(\vect{Q})\,\ee^{\ii\vect{Q}\cdot\vect{R}}\,u_{\vect{Q}\nu} , \qquad
  \vect{p}_{\vect{R}\kappa} = \sqrt{\frac{M_\kappa}{N}}\sum_{\vect{Q}\nu}\vect{e}_{\kappa\nu}(\vect{Q})^*\,\ee^{-\ii\vect{Q}\cdot\vect{R}}\,p_{\vect{Q}\nu} .
  \label{eq:5.3-inverse}
\end{equation}
Inserting \cref{eq:5.3-inverse} into \cref{eq:5.3-normal} returns $u_{\vect{Q}\nu}$, by the lattice sum and the orthonormality of the polarization vectors. The normal coordinates are not Hermitian. Since $\vect{u}_{\vect{R}\kappa}$ is, and $\vect{e}_{\kappa\nu}(-\vect{Q}) = \vect{e}_{\kappa\nu}(\vect{Q})^*$ (\cref{eq:A.6-e}), the adjoint of \cref{eq:5.3-normal} is the coordinate at $-\vect{Q}$, and the same holds for the momenta:
\begin{equation*}
  u_{\vect{Q}\nu}^\dagger = u_{-\vect{Q}\nu} , \qquad p_{\vect{Q}\nu}^\dagger = p_{-\vect{Q}\nu} .
\end{equation*}
The commutator follows from that of \cref{eq:5.3-H}. The masses cancel, the lattice sum sets $\vect{Q}' = \vect{Q}$, and orthonormality sets $\nu' = \nu$:
\begin{equation*}
  [u_{\vect{Q}\nu}, p_{\vect{Q}'\nu'}] = \frac{\ii}{N}\sum_{\vect{R}\kappa}\ee^{\ii(\vect{Q}' - \vect{Q})\cdot\vect{R}}\,\vect{e}_{\kappa\nu}(\vect{Q})^*\cdot\vect{e}_{\kappa\nu'}(\vect{Q}') = \ii\,\delta_{\vect{Q}\vect{Q}'}\delta_{\nu\nu'} ,
\end{equation*}
while the coordinates commute among themselves, and so do the momenta. Each pair $(u_{\vect{Q}\nu}, p_{\vect{Q}\nu})$ is as canonical as the position and momentum of a single particle; the only novelty is that the pairs at $\vect{Q}$ and $-\vect{Q}$ are adjoints of each other. In field theory the momentum conjugate to a mode of wavevector $\vect{k}$ carries $-\vect{k}$; here the phase $\ee^{+\ii\vect{Q}\cdot\vect{R}}$ in the definition of $p_{\vect{Q}\nu}$ absorbs the sign, so that $p_{\vect{Q}\nu}$ is conjugate to $u_{\vect{Q}\nu}$ itself.

In these variables the Hamiltonian separates. In the kinetic term we write $\vect{p}^2_{\vect{R}\kappa} = \vect{p}^\dagger_{\vect{R}\kappa}\cdot\vect{p}_{\vect{R}\kappa}$ and insert \cref{eq:5.3-inverse}; the sum over $\vect{R}$ sets the two wavevectors equal and the sum over $\kappa$ the two branches. In the potential term we write the left displacement as $\vect{u}^\dagger_{\vect{R}\kappa}$ and expand both. The sum over $\vect{R}'$ of the force constants times $\ee^{\ii\vect{Q}'\cdot\vect{R}'}$ is $\ee^{\ii\vect{Q}'\cdot\vect{R}}\sqrt{M_\kappa M_{\kappa'}}\,D^{\mathrm{dyn}}_{\kappa\kappa'}(\vect{Q}')$, by translation invariance and \cref{eq:5.2-Ddyn}. The dynamical matrix acting on its eigenvector returns $\omega^2_{\vect{Q}'\nu'}$, and the sums over $\vect{R}$ and $\kappa$ again set $\vect{Q}' = \vect{Q}$ and $\nu' = \nu$. What remains is
\begin{equation}
  H_{\mathrm{N}} = \frac12\sum_{\vect{Q}\nu}\Bigl(p^\dagger_{\vect{Q}\nu}p_{\vect{Q}\nu} + \omega_{\vect{Q}\nu}^2\,u^\dagger_{\vect{Q}\nu}u_{\vect{Q}\nu}\Bigr) .
  \label{eq:5.3-Hnormal}
\end{equation}
The $3n_{\mathrm{at}}N$ coupled nuclei have become $3n_{\mathrm{at}}N$ independent oscillators of unit mass, one for each mode. With \cref{eq:5.3-Hnormal} the Heisenberg equations are $\dot u_{\vect{Q}\nu} = p^\dagger_{\vect{Q}\nu}$ and $\dot p^\dagger_{\vect{Q}\nu} = -\omega_{\vect{Q}\nu}^2u_{\vect{Q}\nu}$, so that $\ddot u_{\vect{Q}\nu} = -\omega^2_{\vect{Q}\nu}u_{\vect{Q}\nu}$: the classical equation of each normal mode, as \cref{sec:5.2} anticipated. Nothing quantum has happened yet. It enters through the commutator, which forbids $u_{\vect{Q}\nu}$ and $p_{\vect{Q}\nu}$ from vanishing together.

\subsection{Phonon operators}

For a single oscillator of unit mass the ladder operator is $(\omega/2)^{1/2}x + \ii(2\omega)^{-1/2}p$. The normal modes need one change: since $u_{\vect{Q}\nu}$ and $p_{\vect{Q}\nu}$ are not Hermitian, the momentum enters through its adjoint,
\begin{equation}
  b_{\vect{Q}\nu} = \Bigl(\frac{\omega_{\vect{Q}\nu}}{2}\Bigr)^{1/2}u_{\vect{Q}\nu} + \frac{\ii}{(2\omega_{\vect{Q}\nu})^{1/2}}\,p^\dagger_{\vect{Q}\nu} , \qquad
  b^\dagger_{\vect{Q}\nu} = \Bigl(\frac{\omega_{\vect{Q}\nu}}{2}\Bigr)^{1/2}u_{-\vect{Q}\nu} - \frac{\ii}{(2\omega_{\vect{Q}\nu})^{1/2}}\,p_{\vect{Q}\nu} .
  \label{eq:5.3-b}
\end{equation}
With the commutator of the normal coordinates and its adjoint, $[u_{\vect{Q}\nu}, p^\dagger_{\vect{Q}'\nu'}] = \ii\,\delta_{\vect{Q},-\vect{Q}'}\delta_{\nu\nu'}$, the two cross terms of $[b_{\vect{Q}\nu}, b^\dagger_{\vect{Q}'\nu'}]$ contribute $\tfrac12\delta_{\vect{Q}\vect{Q}'}\delta_{\nu\nu'}$ each. The two cross terms of $[b_{\vect{Q}\nu}, b_{\vect{Q}'\nu'}]$ are $\mp\tfrac12\delta_{\vect{Q},-\vect{Q}'}\delta_{\nu\nu'}$ and cancel, because $\omega_{-\vect{Q}\nu} = \omega_{\vect{Q}\nu}$. Hence
\begin{equation*}
  [b_{\vect{Q}\nu}, b^\dagger_{\vect{Q}'\nu'}] = \delta_{\vect{Q}\vect{Q}'}\delta_{\nu\nu'} , \qquad
  [b_{\vect{Q}\nu}, b_{\vect{Q}'\nu'}] = [b^\dagger_{\vect{Q}\nu}, b^\dagger_{\vect{Q}'\nu'}] = 0 ,
\end{equation*}
the boson algebra \cref{eq:2.3-comm} with $\gamma = (\vect{Q}, \nu)$. Adding \cref{eq:5.3-b} at $\vect{Q}$ to the adjoint at $-\vect{Q}$, and subtracting, inverts the definition:
\begin{equation*}
  u_{\vect{Q}\nu} = \frac{b_{\vect{Q}\nu} + b^\dagger_{-\vect{Q}\nu}}{(2\omega_{\vect{Q}\nu})^{1/2}} , \qquad
  p^\dagger_{\vect{Q}\nu} = -\ii\Bigl(\frac{\omega_{\vect{Q}\nu}}{2}\Bigr)^{1/2}\bigl(b_{\vect{Q}\nu} - b^\dagger_{-\vect{Q}\nu}\bigr) .
\end{equation*}
A normal coordinate contains the annihilator of the mode $\vect{Q}$ and the creator of the mode $-\vect{Q}$. Inserting these forms into \cref{eq:5.3-Hnormal}, one bracket is
\begin{equation*}
  p^\dagger_{\vect{Q}\nu}p_{\vect{Q}\nu} + \omega^2_{\vect{Q}\nu}u^\dagger_{\vect{Q}\nu}u_{\vect{Q}\nu}
  = \frac{\omega_{\vect{Q}\nu}}{2}\Bigl[(b_{\vect{Q}\nu} - b^\dagger_{-\vect{Q}\nu})(b^\dagger_{\vect{Q}\nu} - b_{-\vect{Q}\nu}) + (b^\dagger_{\vect{Q}\nu} + b_{-\vect{Q}\nu})(b_{\vect{Q}\nu} + b^\dagger_{-\vect{Q}\nu})\Bigr] .
\end{equation*}
The products of two annihilators, $-b_{\vect{Q}\nu}b_{-\vect{Q}\nu} + b_{-\vect{Q}\nu}b_{\vect{Q}\nu}$, cancel since the two operators commute, and so do the products of two creators. The four remaining terms are $b_{\vect{Q}\nu}b^\dagger_{\vect{Q}\nu} + b^\dagger_{\vect{Q}\nu}b_{\vect{Q}\nu}$ and the same at $-\vect{Q}$, which in the sum over all $\vect{Q}$ repeats the first pair. With $bb^\dagger = b^\dagger b + 1$,
\begin{equation}
  H_{\mathrm{N}} = \sum_{\vect{Q}\nu}\omega_{\vect{Q}\nu}\Bigl(b^\dagger_{\vect{Q}\nu}b_{\vect{Q}\nu} + \frac12\Bigr) .
  \label{eq:5.3-Hb}
\end{equation}

The spectrum is that of independent oscillators. The eigenstates are labeled by one occupation number $n_{\vect{Q}\nu} = 0, 1, 2, \dots$ per mode, $|\{n_{\vect{Q}\nu}\}\rangle = \prod_{\vect{Q}\nu}(n_{\vect{Q}\nu}!)^{-1/2}(b^\dagger_{\vect{Q}\nu})^{n_{\vect{Q}\nu}}|0\rangle$, with energy $\sum_{\vect{Q}\nu}\omega_{\vect{Q}\nu}(n_{\vect{Q}\nu} + \tfrac12)$. A \emph{phonon} is one quantum of a normal mode: the state $b^\dagger_{\vect{Q}\nu}|0\rangle$ lies $\omega_{\vect{Q}\nu}$ above the ground state, and $\omega_{\vect{Q}\nu}$ of \cref{sec:5.2} is its energy. Phonons are bosons, since any number of them can occupy a mode, and their number is not conserved, which is why their thermal occupation has no chemical potential (\cref{sec:2.5}). Each phonon also carries a crystal momentum. A lattice translation $T_{\vect{R}}$, in the sense of \cref{eq:3.1-crystalmomentum}, relabels the displacements, $\vect{u}_{\vect{R}'\kappa} \to \vect{u}_{\vect{R}'-\vect{R},\kappa}$. By \cref{eq:5.3-normal,eq:5.3-b} this multiplies $b^\dagger_{\vect{Q}\nu}$ by $\ee^{\ii\vect{Q}\cdot\vect{R}}$, while the ground state is invariant. A state with one phonon in the mode $\vect{Q}\nu$ has crystal momentum $\vect{Q}$, defined up to a reciprocal-lattice vector, exactly as an electron in a Bloch state has $\vect{k}$ (\cref{sec:3.1}).

The ground state $|0\rangle$ holds no phonons, but its energy is not zero. Every mode keeps $\tfrac12\omega_{\vect{Q}\nu}$, and in silicon the zero-point energy $\tfrac12\int\dd\omega\,g_{\mathrm{ph}}(\omega)\,\omega$ is \SI{123}{meV} per cell for the computed density of states of \cref{fig:5-si}: $\tfrac12\times\SI{40.9}{meV}$, the mean phonon energy, times six modes. It depends on the cell volume through the frequencies, which \cref{sec:5.7} takes up. In the Heisenberg picture the phonon operators of \cref{eq:5.3-Hb} evolve as $b_{\vect{Q}\nu}(t) = b_{\vect{Q}\nu}\ee^{-\ii\omega_{\vect{Q}\nu}t}$, and every displacement is a superposition of oscillations at the mode frequencies, as in the classical crystal.

The displacement of a nucleus follows from \cref{eq:5.3-inverse} once the normal coordinate is written with the phonon operators. For the momentum, the term $-\vect{Q}$ of the sum is relabeled with $\vect{e}_{\kappa\nu}(-\vect{Q})^* = \vect{e}_{\kappa\nu}(\vect{Q})$. Both have the same structure:
\begin{keyresult}[Phonon expansion of the displacement]
\begin{equation}
  \vect{u}_{\vect{R}\kappa} = \sum_{\vect{Q}\nu}\frac{1}{\sqrt{2NM_\kappa\omega_{\vect{Q}\nu}}}\;\vect{e}_{\kappa\nu}(\vect{Q})\,\ee^{\ii\vect{Q}\cdot\vect{R}}\,B_{\vect{Q}\nu} , \qquad
  B_{\vect{Q}\nu} = b_{\vect{Q}\nu} + b^\dagger_{-\vect{Q}\nu} ,
  \label{eq:5.3-u}
\end{equation}
\begin{equation*}
  \vect{p}_{\vect{R}\kappa} = -\ii\sum_{\vect{Q}\nu}\Bigl(\frac{M_\kappa\omega_{\vect{Q}\nu}}{2N}\Bigr)^{1/2}\vect{e}_{\kappa\nu}(\vect{Q})\,\ee^{\ii\vect{Q}\cdot\vect{R}}\,\bigl(b_{\vect{Q}\nu} - b^\dagger_{-\vect{Q}\nu}\bigr) ,
\end{equation*}
with $[b_{\vect{Q}\nu}, b^\dagger_{\vect{Q}'\nu'}] = \delta_{\vect{Q}\vect{Q}'}\delta_{\nu\nu'}$, $\sum_\kappa|\vect{e}_{\kappa\nu}(\vect{Q})|^2 = 1$ and $\vect{e}_{\kappa\nu}(-\vect{Q}) = \vect{e}_{\kappa\nu}(\vect{Q})^*$. The harmonic Hamiltonian is \cref{eq:5.3-Hb}.
\end{keyresult}
This is \cref{eq:A.6-u}, the convention used throughout, and that of Ashcroft and Mermin~\cite[Chap.~23]{ashcroft1976solid} and of Giustino~\cite[Eqs.~(16)--(21)]{giustino2017electron}, with the phonon wavevector written $\vect{Q}$. The \emph{displacement field} $B_{\vect{Q}\nu}$ obeys $B^\dagger_{\vect{Q}\nu} = B_{-\vect{Q}\nu}$, so that the terms at $\vect{Q}$ and $-\vect{Q}$ of \cref{eq:5.3-u} are adjoints of each other and the displacement is Hermitian. It is the operator that every coupling linear in the displacements contains: the electron--phonon Hamiltonian \cref{eq:A.6-Hep} and the single-phonon response of a nucleus to dark matter.

\subsection{The zero-point length}

Each term of \cref{eq:5.3-u} has three factors: the pattern $\vect{e}_{\kappa\nu}(\vect{Q})\ee^{\ii\vect{Q}\cdot\vect{R}}$ of \cref{sec:5.2}, the operator $B_{\vect{Q}\nu}$, which changes the number of phonons by one, and a length that fixes how far one quantum moves the nuclei. We write that length as the \emph{zero-point length} of the mode,
\begin{equation}
  \ell_{\vect{Q}\nu} = \frac{1}{\sqrt{2M_{\mathrm{cell}}\,\omega_{\vect{Q}\nu}}} ,
  \label{eq:5.3-ell}
\end{equation}
with the cell mass $M_{\mathrm{cell}} = \sum_\kappa M_\kappa$ and the factor $1/\sqrt N$ kept outside (\cref{eq:A.6-ell}). For an acoustic mode at small $\vect{Q}$, where $\vect{e}_{\kappa\nu} = (M_\kappa/M_{\mathrm{cell}})^{1/2}\hat{\vect{e}}_\nu$ (\cref{sec:5.2}), the amplitude in \cref{eq:5.3-u} is $\ell_{\vect{Q}\nu}/\sqrt N$ for every nucleus: the whole cell moves as one particle of mass $M_{\mathrm{cell}}$. For other modes $\ell_{\vect{Q}\nu}$ sets the scale, and the polarization vector distributes the motion over the nuclei of the cell. The name refers to the ground state, in which $\langle0|B_{\vect{Q}\nu}B^\dagger_{\vect{Q}\nu}|0\rangle = \langle0|b_{\vect{Q}\nu}b^\dagger_{\vect{Q}\nu}|0\rangle = 1$: $\ell_{\vect{Q}\nu}/\sqrt N$ is the zero-point amplitude that one mode contributes to each nucleus.

For silicon, with $M_{\mathrm{cell}}$ two silicon masses (\SI{52.33}{GeV}, \cref{sec:1.5}), the zone-center optical phonon at \SI{64.1}{meV} has $\ell_{\vect{Q}\nu} = \SI{0.024}{\angstrom}$. This is half the zero-point length $\ell_{\mathrm{zp}} = \SI{0.048}{\angstrom}$ of the bond coordinate in \cref{fig:5-scales}(c), and must be. That coordinate is the change of a bond length, and in the optical mode the two atoms of the cell move against each other with equal amplitude, $\vect{e}_{\kappa\nu} = \pm\hat{\vect{e}}/\sqrt2$, each by half of it. The transverse acoustic phonon at X, \SI{18.6}{meV}, has $\ell_{\vect{Q}\nu} = \SI{0.045}{\angstrom}$, and an acoustic phonon of \SI{5}{meV} has \SI{0.086}{\angstrom}: soft modes move the nuclei farther, as $\omega_{\vect{Q}\nu}^{-1/2}$. The factor $1/\sqrt N$ makes the motion of any single mode unobservably small in a macroscopic crystal. A kilogram of silicon has $N = \num{1.07e25}$ cells (\cref{sec:3.1}), and one optical phonon moves each of its nuclei by $\SI{7e-15}{\angstrom}$. What is finite is the sum over all $3n_{\mathrm{at}}N$ modes, each contributing a share of order $\ell^2_{\vect{Q}\nu}/N$ to the mean-square displacement, and we compute that sum below.

\subsection{Emission and absorption}

The matrix elements of the displacement between phonon states follow from \cref{eq:5.3-u} and the boson relations \cref{eq:2.3-boson}. Only two terms of the sum connect a state with $n_{\vect{Q}\nu}$ phonons in the mode $\vect{Q}\nu$ to a state that differs from it in that mode alone. The term $\vect{Q}$, through $b_{\vect{Q}\nu}$, removes a phonon. The term $-\vect{Q}$, through $b^\dagger_{\vect{Q}\nu}$, adds one, with the pattern $\vect{e}_{\kappa\nu}(-\vect{Q})\ee^{-\ii\vect{Q}\cdot\vect{R}} = \vect{e}_{\kappa\nu}(\vect{Q})^*\ee^{-\ii\vect{Q}\cdot\vect{R}}$. Hence
\begin{equation}
  \begin{gathered}
  \langle n_{\vect{Q}\nu} + 1|\,\vect{u}_{\vect{R}\kappa}\,|n_{\vect{Q}\nu}\rangle = \frac{\vect{e}_{\kappa\nu}(\vect{Q})^*\,\ee^{-\ii\vect{Q}\cdot\vect{R}}}{\sqrt{2NM_\kappa\omega_{\vect{Q}\nu}}}\,\sqrt{n_{\vect{Q}\nu} + 1} ,
  \\
  \langle n_{\vect{Q}\nu} - 1|\,\vect{u}_{\vect{R}\kappa}\,|n_{\vect{Q}\nu}\rangle = \frac{\vect{e}_{\kappa\nu}(\vect{Q})\,\ee^{\ii\vect{Q}\cdot\vect{R}}}{\sqrt{2NM_\kappa\omega_{\vect{Q}\nu}}}\,\sqrt{n_{\vect{Q}\nu}} ,
  \end{gathered}
  \label{eq:5.3-me}
\end{equation}
with the occupations of all other modes unchanged. In thermal equilibrium the squared matrix elements are averaged over the Gibbs state of \cref{sec:2.5}, and the occupation becomes the mean \emph{phonon occupation}
\begin{equation*}
  n_{\vect{Q}\nu} = \langle b^\dagger_{\vect{Q}\nu}b_{\vect{Q}\nu}\rangle = n_{\mathrm{B}}(\omega_{\vect{Q}\nu}) ,
\end{equation*}
by \cref{eq:2.5-nF}. Emission then carries the weight $n_{\vect{Q}\nu} + 1$ and absorption the weight $n_{\vect{Q}\nu}$, the factors $n_{\vect{Q}\nu} + \tfrac12 \pm \tfrac12$ of \cref{eq:A.6-weights}, and their ratio $\ee^{\beta\omega_{\vect{Q}\nu}}$ is detailed balance for the mode. Each mode of the crystal is the oscillator of \cref{we:2-oscillator}, with the coupling $\lambda_{\mathrm{osc}}$ replaced by the projection of the dark-matter vertex on $\vect{e}_{\kappa\nu}(\vect{Q})/\sqrt{2NM_\kappa\omega_{\vect{Q}\nu}}$. The numbers of that example carry over. At room temperature the optical phonon of silicon has $n_{\vect{Q}\nu} = 0.09$ and the transverse acoustic phonon at X has $0.95$. In a cryogenic detector every mode above a few tens of \si{\micro eV} is empty, absorption is absent, and emission into an empty mode, with weight one, is the whole signal.

Consider an interaction of dark matter with the nuclei of one species $\kappa$ that depends on their positions through $\ee^{\ii\vect{q}\cdot\vect{X}_I}$, the nuclear density $n_\kappa(\vect{q})$ of \cref{sec:2.3}, with $\vect{q}$ the momentum transfer. With $\vect{X} = \vect{R} + \vect{\tau}_\kappa + \vect{u}_{\vect{R}\kappa}$ (\cref{eq:5.1-X}),
\begin{equation*}
  n_\kappa(\vect{q}) = \sum_{\vect{R}}\ee^{\ii\vect{q}\cdot(\vect{R} + \vect{\tau}_\kappa)}\,\ee^{\ii\vect{q}\cdot\vect{u}_{\vect{R}\kappa}}
  \simeq \sum_{\vect{R}}\ee^{\ii\vect{q}\cdot(\vect{R} + \vect{\tau}_\kappa)}\bigl(1 + \ii\vect{q}\cdot\vect{u}_{\vect{R}\kappa}\bigr) .
\end{equation*}
The expansion holds as long as $\langle(\vect{q}\cdot\vect{u}_{\vect{R}\kappa})^2\rangle = q^2\langle u_x^2\rangle \ll 1$, that is, for $q$ well below \SI{26}{keV} in silicon at room temperature (with the value of $\langle u_x^2\rangle$ given below); the full exponential is the subject of \cref{sec:5.6}. The term 1 leaves the lattice unchanged. The term linear in $\vect{u}$, with \cref{eq:5.3-me} and the lattice sum \cref{eq:3.1-latticesum}, emits one phonon with the amplitude
\begin{equation*}
  \langle n_{\vect{Q}\nu} + 1|\,n_\kappa(\vect{q})\,|n_{\vect{Q}\nu}\rangle
  = \ii\,\ee^{\ii\vect{q}\cdot\vect{\tau}_\kappa}\,\frac{\vect{q}\cdot\vect{e}_{\kappa\nu}(\vect{Q})^*}{\sqrt{2M_\kappa\omega_{\vect{Q}\nu}}}\,\sqrt{N(n_{\vect{Q}\nu} + 1)}\;\sum_{\vect{G}}\delta_{\vect{q},\vect{Q} + \vect{G}} .
\end{equation*}
The lattice sum conserves crystal momentum, a feature of every later phonon rate: the emitted phonon carries the momentum transfer, up to a reciprocal-lattice vector, so that each $\vect{q}$ excites one wavevector $\vect{Q}$ in each branch. The squared amplitude is also proportional to $N$, as it must be for a target whose rate grows with its size. Per nucleus, the weight of the branch $\nu$ is $|\vect{q}\cdot\vect{e}_{\kappa\nu}(\vect{Q})|^2(n_{\vect{Q}\nu} + 1)/(2M_\kappa\omega_{\vect{Q}\nu})$: the projection of the momentum transfer on the polarization, which favors longitudinal modes at small $\vect{Q}$, divided by the energy of the mode, which favors soft ones. Absorption of a phonon $-\vect{q}$ follows in the same way with $n_{\vect{Q}\nu}$, and with several species the amplitudes of the nuclei of a cell add with their phases $\ee^{\ii\vect{q}\cdot\vect{\tau}_\kappa}$.

\subsection{Zero-point and thermal motion}

Each Cartesian component of a nuclear displacement behaves roughly like an oscillator of mass $M_\kappa$ at a typical phonon frequency $\bar\omega$. In the ground state such an oscillator has $\langle x^2\rangle = 1/(2M_\kappa\bar\omega)$, and at a temperature well above $\bar\omega$ equipartition, $M_\kappa\bar\omega^2\langle x^2\rangle = T$, leads to $T/(M_\kappa\bar\omega^2)$. For silicon with $\bar\omega \sim \SI{30}{meV}$ the zero-point estimate is $\langle u_x^2\rangle^{1/2} \sim \SI{0.05}{\angstrom}$, the number of \cref{we:5-kappa}, and room temperature, where $T = \SI{25.9}{meV}$ is close to $\bar\omega$, should roughly double $\langle u_x^2\rangle$. The question is what $\bar\omega$ means for a crystal with a whole spectrum of modes.

The exact statement is a sum over modes. In the Gibbs state the occupations are independent, and the only nonzero averages of two displacement fields pair a mode with its partner at $-\vect{Q}$:
\begin{equation*}
  \langle B_{\vect{Q}\nu}B_{\vect{Q}'\nu'}\rangle = \delta_{\vect{Q}',-\vect{Q}}\delta_{\nu\nu'}\bigl(\langle b_{\vect{Q}\nu}b^\dagger_{\vect{Q}\nu}\rangle + \langle b^\dagger_{-\vect{Q}\nu}b_{-\vect{Q}\nu}\rangle\bigr) = \delta_{\vect{Q}',-\vect{Q}}\delta_{\nu\nu'}\,(2n_{\vect{Q}\nu} + 1) ,
\end{equation*}
by \cref{eq:2.5-nF} and $n_{-\vect{Q}\nu} = n_{\vect{Q}\nu}$. Inserting \cref{eq:5.3-u} twice, the phases $\ee^{\ii\vect{Q}\cdot\vect{R}}\ee^{-\ii\vect{Q}\cdot\vect{R}}$ cancel, $\vect{e}_{\kappa\nu}(-\vect{Q}) = \vect{e}_{\kappa\nu}(\vect{Q})^*$ supplies the conjugate, and the \emph{displacement tensor} of the nucleus $\kappa$ is
\begin{equation}
  B^{\alpha\beta}_\kappa = \langle u^\alpha_{\vect{R}\kappa}u^\beta_{\vect{R}\kappa}\rangle
  = \frac1N\sum_{\vect{Q}\nu}\frac{2n_{\vect{Q}\nu} + 1}{2M_\kappa\,\omega_{\vect{Q}\nu}}\;e^\alpha_{\kappa\nu}(\vect{Q})\,e^\beta_{\kappa\nu}(\vect{Q})^* ,
  \qquad 2n_{\vect{Q}\nu} + 1 = \coth\frac{\omega_{\vect{Q}\nu}}{2T} ,
  \label{eq:5.3-B}
\end{equation}
the same in every cell. The terms at $\vect{Q}$ and $-\vect{Q}$ are complex conjugates of each other, which makes $B_\kappa$ real and symmetric, and positive since every term of the diagonal is. The average is over the thermal state, and at $T = 0$ over the ground state: $B_\kappa$ then measures the \emph{zero-point motion}, the spread of the nuclear wavefunction that no cooling removes. The two limits are those of the estimate. At $T = 0$, $2n + 1 = 1$, and each mode contributes $1/(2M_\kappa\omega_{\vect{Q}\nu})$ times the weight $|\vect{e}_{\kappa\nu}|^2/N$. For $T$ above every phonon energy, $2n + 1 \to 2T/\omega_{\vect{Q}\nu}$, and each mode contributes its equipartition value $T/(M_\kappa\omega^2_{\vect{Q}\nu})$, growing linearly in $T$. The acoustic modes have $\omega_{\vect{Q}\nu} \to 0$, but they are few: in three dimensions their number per unit frequency falls as $\omega^2$ (\cref{eq:5.2-dosac}), and the sum stays finite at every temperature.

At a site of cubic symmetry, such as the tetrahedral sites of the diamond and zinc-blende structures, a symmetric second-rank tensor that is invariant under the site group is a multiple of the identity. There $B_\kappa = \langle u_x^2\rangle\,\mathbf{1}$, and the mean-square displacement $\langle u^2\rangle = \operatorname{tr}B_\kappa = 3\langle u_x^2\rangle$; crystallographers quote $\langle u_x^2\rangle$, or the isotropic displacement parameter $8\pi^2\langle u_x^2\rangle$. If, in addition, the $n_{\mathrm{at}}$ nuclei of the cell are of one species and related by symmetry, as in silicon and germanium, all of them have the same trace. We can then average the trace of \cref{eq:5.3-B} over $\kappa$, the normalization $\sum_\kappa|\vect{e}_{\kappa\nu}|^2 = 1$ removes the polarization vectors, and the sum over modes becomes an integral over the density of states \cref{eq:5.2-dos}:
\begin{equation}
  \langle u^2\rangle = \int_0^\infty\dd\omega\;\frac{g_{\mathrm{ph}}(\omega)}{n_{\mathrm{at}}}\;\frac{\coth(\omega/2T)}{2M_\kappa\,\omega} .
  \label{eq:5.3-u2}
\end{equation}
The density of states enters per atom, $g_{\mathrm{ph}}/n_{\mathrm{at}}$, normalized to three. For a crystal of several species, such as GaAs, the same steps keep the weight of each nucleus separately, and $g_{\mathrm{ph}}/n_{\mathrm{at}}$ is replaced by the atom-projected density of states $g_{\mathrm{ph},\kappa}(\omega) = N^{-1}\sum_{\vect{Q}\nu}|\vect{e}_{\kappa\nu}(\vect{Q})|^2\delta(\omega - \omega_{\vect{Q}\nu})$, also normalized to three. \Cref{eq:5.3-u2} makes the estimate precise: the typical frequency at $T = 0$ is the inverse of the mean of $1/\omega$ over the spectrum, $\langle u^2\rangle = (3/2M_\kappa)\langle\omega^{-1}\rangle$, and at high temperature the inverse root of the mean of $1/\omega^2$, $\langle u^2\rangle = (3T/M_\kappa)\langle\omega^{-2}\rangle$. Both means are weighted toward low frequencies, the second more strongly.

In the Debye model, $g_{\mathrm{ph}}/n_{\mathrm{at}} = 9\omega^2/\omega_{\mathrm{D}}^3$ below $\omega_{\mathrm{D}}$ (\cref{eq:5.2-gD}). Splitting $\coth = 1 + 2n_{\mathrm{B}}$ into its zero-point and thermal parts and writing $x = \omega/T$,
\begin{equation}
  \langle u^2\rangle = \frac{9}{4M_\kappa\omega_{\mathrm{D}}}\Bigl[1 + 4\Bigl(\frac{T}{\omega_{\mathrm{D}}}\Bigr)^2\int_0^{\omega_{\mathrm{D}}/T}\frac{x\,\dd x}{\ee^x - 1}\Bigr] ,
  \label{eq:5.3-debye}
\end{equation}
from which the Debye--Waller exponent \cref{eq:D.15-debye} follows. The zero-point value $9/(4M_\kappa\omega_{\mathrm{D}})$ corresponds to $\bar\omega = 2\omega_{\mathrm{D}}/3$ in the estimate. At high temperature the integral is $y - y^2/4 + y^3/36 - \cdots$ with $y = \omega_{\mathrm{D}}/T$. The bracket then becomes $4T/\omega_{\mathrm{D}} + \omega_{\mathrm{D}}/(9T) + \cdots$, and $\langle u^2\rangle = (9T/M_\kappa\omega_{\mathrm{D}}^2)[1 + \omega_{\mathrm{D}}^2/36T^2 + \cdots]$, the equipartition value with the zero-point motion as its first correction. The thermal part equals the zero-point part when the bracket is two, at $T = 0.44\,\Theta_{\mathrm{D}}$. For silicon this is \SI{284}{K}: at room temperature the nuclei of silicon move as much because of their zero-point motion as because of the heat.

For crystals of similar bonding, the zero-point value of \cref{eq:5.3-debye} $\propto 1/(M_\kappa\omega_{\mathrm{D}})$ falls with the mass, since $\omega_{\mathrm{D}} \propto M_\kappa^{-1/2}$ at fixed force constants: heavier nuclei are better localized. The high-temperature value $\propto T/(M_\kappa\omega_{\mathrm{D}}^2)$ depends on the force constants alone, the classical statement that a spring stretched by thermal energy does not care what mass hangs on it. A heavier and softer crystal has less zero-point motion and, at room temperature, more thermal motion, as germanium shows in \cref{tab:5-phonons}.

In practice $B_\kappa$ is computed from DFPT phonons. The sum \cref{eq:5.3-B} runs over a uniform grid of wavevectors, with the frequencies and polarization vectors interpolated from the force constants of a coarser grid, or \cref{eq:5.3-u2} is integrated over a density of states computed on a fine grid. The grid misses the small acoustic frequencies near $\mathrm{\Gamma}$, whose terms grow as $1/\omega$ at $T = 0$ and $1/\omega^2$ at $T > 0$, and \cref{eq:E.4-dBgamma} adds them back analytically. For silicon the $4^3$ grid of \cref{fig:5-scales} yields $\langle u_x^2\rangle^{1/2} = \SI{0.050}{\angstrom}$ at zero temperature and \SI{0.081}{\angstrom} at \SI{293}{K}, and the $48^3$ density of states of \cref{fig:5-si} leads to \SI{0.0495}{\angstrom} at zero temperature, \SI{0.0799}{\angstrom} at \SI{293}{K} and \SI{0.0807}{\angstrom} at \SI{300}{K}. The two grids differ by \SI{2}{\percent} at most, well inside the error that the functional itself makes (\cref{we:5-u2}).

\begin{workedexample}{Zero-point and room-temperature motion of silicon}\label{we:5-u2}
We evaluate \cref{eq:5.3-debye} for silicon with the natural isotope mixture, $M_\kappa = \SI{28.0855}{u} = \SI{26.16}{GeV}$. The Debye frequency is that of the measured elastic constants, $\omega_{\mathrm{D}} = \SI{55.6}{meV}$ (\cref{we:5-debye}). We convert with $\hbar c = \SI{1973.27}{eV.\angstrom}$.

\emph{Zero temperature.} The product $M_\kappa\omega_{\mathrm{D}} = \SI{1.455e9}{eV^2}$, and
\begin{equation*}
  \langle u^2\rangle_{T=0} = \frac{9}{4M_\kappa\omega_{\mathrm{D}}} = \SI{1.546e-9}{eV^{-2}} = \SI{6.02e-3}{\angstrom^2} , \qquad
  \langle u_x^2\rangle^{1/2}_{T=0} = \SI{0.0448}{\angstrom} .
\end{equation*}

\emph{Room temperature.} At $T = \SI{300}{K} = \SI{25.85}{meV}$, $\omega_{\mathrm{D}}/T = 2.152$, and the integral of \cref{eq:5.3-debye}, evaluated numerically, is $\int_0^{2.152}x\,\dd x/(\ee^x - 1) = 1.259$. The bracket is $1 + 4\times0.4648^2\times1.259 = 2.088$, so that
\begin{equation*}
  \langle u^2\rangle_{\SI{300}{K}} = \SI{1.257e-2}{\angstrom^2} , \qquad
  \langle u_x^2\rangle^{1/2}_{\SI{300}{K}} = \SI{0.0647}{\angstrom} .
\end{equation*}
Thermal phonons double the mean-square displacement, as the estimate suggested, and multiply the amplitude by $\sqrt{2.088} = 1.44$. Since $T$ lies below $\omega_{\mathrm{D}}/2$, equipartition alone, $9T/(M_\kappa\omega^2_{\mathrm{D}})$, falls \SI{11}{\percent} short, at $\langle u_x^2\rangle = \SI{3.73e-3}{\angstrom^2}$; its first zero-point correction, the factor $1 + \omega_{\mathrm{D}}^2/36T^2 = 1.129$, recovers the full value to \SI{0.5}{\percent}.

\emph{Against the full spectrum and the measurement.} The value at \SI{293}{K} from a lattice-dynamical fit to the measured phonons, the reference of \cref{we:5-kappa}, is $\langle u_x^2\rangle = \SI{5.941e-3}{\angstrom^2}$~\cite{flensburg1999lattice}. The Debye model predicts \SI{4.11e-3}{\angstrom^2} at that temperature. \Cref{eq:5.3-u2} with the computed density of states of \cref{fig:5-si} gives $\langle u_x^2\rangle = \SI{2.45e-3}{\angstrom^2}$ at zero temperature and \SI{6.38e-3}{\angstrom^2} at \SI{293}{K}, \SI{7}{\percent} above the measurement, as expected from the transverse acoustic branches, which are \SI{8}{\percent} soft in this calculation (\cref{sec:5.2}) and enter as $1/\omega^2$. Against this spectrum the Debye model is low by \SI{18}{\percent} at zero temperature and by about a third at room temperature, and the reason is the flat transverse acoustic branches. In the computed spectrum the modes below \SI{20}{meV} are \SI{22}{\percent} of all modes but carry \SI{45}{\percent} of the zero-point $\langle u^2\rangle$, since they enter with the weight $1/\omega$. In the Debye spectrum the same window holds \SI{5}{\percent} of the modes and \SI{13}{\percent} of $\langle u^2\rangle$. The typical frequencies of the estimate make the same point. For the computed spectrum, $\langle\omega^{-1}\rangle^{-1} = \SI{30.4}{meV}$ and $\langle\omega^{-2}\rangle^{-1/2} = \SI{25.2}{meV}$, against $2\omega_{\mathrm{D}}/3 = \SI{37.1}{meV}$ and $\omega_{\mathrm{D}}/\sqrt3 = \SI{32.1}{meV}$ in the model. Elastic constants fix the low-frequency density of states exactly, but the displacement is weighted toward the zone-boundary modes of the transverse branches, which the model pushes up to $\omega_{\mathrm{D}}$.

The verdict is that the Debye model reproduces the zero-point and room-temperature motion of silicon to \SIrange{20}{35}{\percent}, which is enough for an estimate of where the Debye--Waller suppression sets in. A rate that depends on $B_\kappa$ at the ten-percent level needs the full density of states, and at that level the softness of the computed transverse branches matters as well.
\end{workedexample}

\Cref{tab:5-phonons} puts silicon next to the two contrast materials of the book. Germanium has the same structure and bonding, a nucleus 2.6 times heavier and weaker bonds. Its optical phonon, at \SI{37.6}{meV}, lies below the \SI{39.9}{meV} that the mass alone would give, \SI{64.1}{meV} scaled by the square root of the mass ratio, and its elastic constants imply $\Theta_{\mathrm{D}} = \SI{373}{K}$. Its zero-point motion is smaller than that of silicon and its room-temperature motion larger, as the trends of \cref{eq:5.3-debye} require. At \SI{300}{K}, $T = 0.80\,\Theta_{\mathrm{D}}$, and the thermal part is 2.4 times the zero-point part. Its transverse acoustic branches are flat near the zone boundary as well, at \SI{9.9}{meV} at X, a third of $\omega_{\mathrm{D}}$ as in silicon, and the Debye values in the table should be read as low by a similar fraction. GaAs, with nuclei of nearly equal mass, sits close to germanium, but it is polar: its longitudinal optical phonon at $\mathrm{\Gamma}$ lies \SI{2.7}{meV} above the transverse one, the LO--TO splitting of \cref{sec:5.5}.

\begin{table}[tbp]
  \centering
  \small
  \caption{Phonon energies and root-mean-square displacements $\langle u_x^2\rangle^{1/2}$ along one Cartesian axis, in \AA, of silicon, germanium and GaAs. Phonon energies: measured by inelastic neutron scattering, silicon at room temperature~\cite{dolling1963inelastic,nilsson1972study}, germanium at \SI{80}{K}~\cite{nilsson1971phonon} as tabulated in Ref.~\cite[Table~IV]{haxhijaj2026lattice}, and GaAs~\cite{strauch1990phonon} as tabulated in Ref.~\cite[Table~I]{dehghani2025influence}; for GaAs the transverse (TO) and longitudinal (LO) optical energies at $\mathrm{\Gamma}$. Debye temperatures from the room-temperature elastic constants and lattice constants with the angular average of \cref{eq:5.2-cbar}: silicon from \cref{we:5-debye}, germanium from $C_{11} = 131$, $C_{12} = 49$, $C_{44} = \SI{68}{GPa}$~\cite{mcskimin1963elastic} (as tabulated in Ref.~\cite[Table~III]{haxhijaj2026lattice}) and the room-temperature $a = \SI{5.658}{\angstrom}$ of \cref{tab:3-materials}. Debye columns: \cref{eq:5.3-debye} with the natural isotope masses. DFPT columns: \cref{eq:5.3-u2} at 0 and \SI{300}{K} with the PBE density of states of \cref{fig:5-si} for silicon, on a $48^3$ grid of wavevectors. Germanium has no DFPT entry because PBE leaves it nearly metallic at the measured lattice constant, with a direct gap of \SI{0.05}{eV} and the wrong band order, and metallic at its own. Entries marked $\dagger$ are computed for GaAs with the same method (PBE, $a = \SI{5.638}{\angstrom}$, a $4^3$ grid of wavevectors and the dipole term of \cref{sec:5.5}): the transverse acoustic energy at X, and the displacements of gallium (first) and arsenic (second) from the atom-projected densities of states $g_{\mathrm{ph},\kappa}$ on a $32^3$ grid. The germanium energies are tabulated to \SI{0.1}{THz}, \SI{0.4}{meV}. A dash marks a value not given here.}
  \label{tab:5-phonons}
  \begin{tabular}{@{}lccccccc@{}}
    \toprule
    & \multicolumn{2}{c}{Phonon energy (meV)} & & \multicolumn{2}{c}{Debye} & \multicolumn{2}{c}{DFPT} \\
    \cmidrule(lr){2-3}\cmidrule(lr){5-6}\cmidrule(lr){7-8}
    & optical at $\mathrm{\Gamma}$ & TA at X & $\Theta_{\mathrm{D}}$ (K) & \SI{0}{K} & \SI{300}{K} & \SI{0}{K} & \SI{300}{K} \\
    \midrule
    Si & 64.1 & 18.6 & 645 & 0.045 & 0.065 & 0.049 & 0.081 \\
    Ge & 37.6 & 9.9 & 373 & 0.037 & 0.067 & -- & -- \\
    GaAs & 33.6, 36.4 & 10.5$^\dagger$ & -- & -- & -- & 0.040, 0.039$^\dagger$ & 0.080, 0.078$^\dagger$ \\
    \bottomrule
  \end{tabular}
\end{table}
